
\section{Design} \label{sec:design}

\subsection{The Surface Language} \label{subsec:the-surface-language}
<<<<<<< HEAD
=======
\begin{enumerate}
    \item $\pathvar$ is our path, or ``context.'' It keeps track of which labels are taking on which values currently.
    
    \item $\env$ is our environment. It functions as a dictionary that maps labels to values. Since a single variable may take on multiple values in different contexts, the environment really maps labels to dictionaries which map contexts to values.
    
    \item $\deps$ is our dependency map. Since certain variable are only defined in the context of other variables, we need a map to tell us which other variable must be defined to bring another variable into scope. So the dependency maps labels to lists of labels for this purpose. For instance, consider the symbolic expression $a -> (b -> b)*_{a}$. Since $b$ is defined in the star expression labelled ``$a$'' then $b$ is only defined if $a$ is defined in the current context of an expression. Thus, $\envget{\deps}{b}=[a]$. When using the environment, we assume the dependency map is precomputed. We also assume the dependency map is consulted before looking up a variables specific value in the environment.
    
    \item $\envput{E}{\lab}{x}$ means ``Insert the key-value pair $(\lab, x)$ into $E$.'' if $E$ is $\env$ then we assume that this operation takes the current context and dependency map into account
    
    \item $\envget{E}{\lab}$ means ``get the value for key $x$ in $E$.'' If $E$ is $\env$ then we assume that this operation takes the current context and dependency map into account.
    
    \item $\envhas{E}{\lab}$ means ``check if $\lab$ is in the key set of $E$''
    
    \item $\hastype{x}{T}$ means ``$x$ is of type $T$'' where $T$ is either a base type $\base$, an arbitrary type $\complex$, or. A base type is a type that is not a list, an arbitrary type either a base type or a list of some arbitrary type. TODO, this isn't quite true I don't think... it really means ``surface expression $x$ will generate values of type $T$.'' maybe need new symbol for this?
    
    \item $e \gens x$ means ``symbolic expression $e$ generates value $x$''

    \item $x \translate y$ means ``surface expression $x$ translates to  core expression $y$''
    
    \item $xs ++ ys$ means ``append the list $ys$ to the list $xs$''
    
    \item $x :: xs$ means ``prepend the element $x$ to the list $xs$''
    
    \item $\fresh{x}$ means that value $x$ is a random value that will never be used again in any other expression
    
    \item $\arbitrary{x}$ means that value $x$ is a random value that may be used again in subsequent expressions.
    
    \item $\gen$ is a function that generates and arbitrary value of the correct type (context dependent) and returns it.
\end{enumerate}
>>>>>>> Evan

\begin{figure}
\begin{center}
    \begin{align*}
       \surface &\to [\surface ]                     & \text{list} \\
            &\OR\surface\; '|'\;\surface              & \text{union}    \\
            &\OR\surface\; '|\{\lab\}'\;\surface      & \text{labeled union}    \\
            &\OR\surface*                                     & \text{kleene closure}    \\
            &\OR\surface*\{\lab\}                             & \text{labeled kleene closure}    \\
            &\OR\surface\;\surface                   & \text{sequencing}    \\
            &\OR \lab \asgnarrow\surface                      & \text{binding}    \\
            &\OR (\lab \Longleftarrow b) \asgnarrow\surface   & \text{constrained binding}    \\
            &\OR \lab                                                  & \text{label}    \\
            &\OR .                                                     & \text{any value}    \\
            &\OR (. \Longleftarrow b)                                  & \text{constrained value}    \\
            &\OR <x>                                                   & \text{constructor application}    \\ 
            &\OR \nil                                                  & \text{empty list}    \\
    \end{align*}
    \caption{The Surface Language as a Context Free Grammar}
    \label{fig:the-surface-language}
\end{center}
\end{figure}

While convenient for generating and checking values,
the core language is quite small and awkward
to write examples in.
Hence, we define a larger, more feature rich surface language
which we then translate into the core language.
Figure~\ref{fig:the-surface-language}
shows the grammar of this core language.

\mypara{Features}
Many features of the core language closely resemble the core language--
we have labels, as well as labeled unions, labeled kleene closures, and labeled bindings, for instance.
However, we introduce a number of additional features.
We also have unlabeled variants on unions, closures, and bindings, in the form of the
unlabeled union, unlabled kleene closure, and unlabeled dot ($.$) (which represents any value.)
Unlike in the core language,
labels in the surface language may \textit{optionally} include a predicate,
but are now required to have one.
Finally,
and perhaps most notably,
whitespace based sequencing
(evocative of traditional regular expressions)
to represent lists,
rather than the explicit \cons and \concat of the core language.

\mypara{Translation}
The surface-to-core translation is split into two phases.
The first phase rewrites the surface language to a subset of the surface language,
in which we have only
labeled features (no unlabeled features),
and all labels have associated predicates.
The second phase rewrites this subset of the surface language to the core language.

\textit{Phase 1}
We define the first translation phase as a relation $\translateOne$,
where $\surface_1 \translate \surface_1$ denotes
that the expression $\surface$ is translated to the core language expression $\surface_2$. 
Figure~\ref{fig:translation-rules-phase1}
shows \bill{a selection of?} the rules
used in the phase 1 translation.
% Besides the difference in the representation of lists,
% many features of the core language closely resemble the core language--
% we have labels, as well as labeled unions, labeled kleene closures, and labeled bindings, for instance.
% We also have unlabeled variants on each of these, in the form of the
% unlabeled union, unlabled kleene closure, and unlabeled dot ($.$) (which represents any value.)
The dot, and unlabeled variants of union and kleene closure,
are essentially syntactic sugar for the labeled variant of each construct,
and are translated as such in phase 1.
For instance,
as shown by the rule
\DotNoPredTrans.
the dot operator $.$,
is translated into
introducing and then matching on a fresh unconstrained label, $(\lab' \Longleftarrow True) \mapsto \lab'$.
Similarly, fresh labels are introduced and bound to each unlabeled union and kleene closure
by rules \UnlabelledUnionTrans and \UnlabelledStarBaseTrans, respectively.

Phase 1 also translates
labels that do not have predicates into
labels with predicates
via the \AssignmentNoPredTrans and \DotNoPredTrans rules.
It does this by simply setting such predicates to $True$,
which is trivially satisfied by any value.
\bill{Not technically a predicate}

\mypara{Phase 2}
Phase 2 translates the surface language to the core language,
$\surface \translate \core$.
In doing so,
it resolves arguably the largest semantic deviation between the surface and core language:
the two languages representations of lists.
In the core language, a single element is concatenated to a list
by \cons, and two lists may be concatenated into a single list by \concat.
The surface language does away with these operators,
using whitespace to denote sequencing--
\texttt{5 6 7} would represent the list $[5, 6, 7]$, for example.
Square brackets are used to represent nested lists--
for instance \texttt{[5 6] [7] [8 9]} represents the nested list $[5, 6], [7], [8, 9]$.
\bill{Maybe elaborate on motivation/design choices somewhere}

The rules \ConcatBaseTrans, \ConcatConsTrans,\ConcatSnocTrans, and \ConcatTrans implement this.
They rely on examining the type of the left and right expressions $\surface_1\;\surface_2$ in a sequence.
In the case that both $\surface_1$ and $\surface_2$
 both have type $[\type]$,
\ConcatTrans concatenates the two lists into a single list with type $[\type]$.
If $\surface_1$ has type \type and $\surface_2$ has type $[\type]$,
then  \ConcatConsTrans says that $\surface_1$ is cons-ed onto the list $\surface_2$.
Similarly, if $\surface_1$ has type $[\type]$,
and $\surface_2$ has type $\type$,
then \ConcatSnocTrans snoc's $\surface_2$ on the end of $\surface_1$.
Finally, \ConcatBaseTrans
handles when \textit{both} $\surface_1$ and $\surface_2$ are the same non-list type \base
by combining them into the list $\surface_1 \cons [\surface_2]$.
Any other combination of types
(for example, if $\surface_1$ has type \type and $\surface_2$
has type $[[\type]]$) is rejected with a compile-time \bill{I think? Evan, is it compile time?} error message.
\bill{More motivation of design choices?}

This difference in the representation of lists does have some slight impact 
on the translation of the kleene star, in particular.
In the core language, we require that the type of the repeated
expression be a list.
In the surface language,
we allow applying the kleene star to either a list \textit{or} a base type,
i.e. $5*$ and $[5]*$ would both be lists of 0 or more 5s.
To support this, we have two rules to support translation of
kleene stars,
\StarTrans and \StarBaseTrans.
\StarTrans translates stars $x*\{\lab\}$ where $x$ is a list type,
by simply translating $x \translate x'$ and then returning $x'*_{\lab}$.
\StarBaseTrans translates stars $x*\{\lab\}$ where $x$ is a base type,
and works by translating of $x$ to $x'$, and then wrapping $x'$ in a list in the star,
$[x']*_{\lab}$.



\bill{...}



\par \bill{Incorporate} The surface language is the language that the user interacts with.
Their expressions are defined in terms of the surface language so they can be translated to the core language.
In this specification, $\lab$ is a label (i.e a string), $x$ is any arbitrary expression, and $b$ is some expression that evaluates to a boolean.
The $b$ expression exists so that the user can constrain the domains of variables so that the generation process isn't completely random. If $b$ is just ``$True$'' then SymCheck will use the first value it generates for the label.
It is possible, however, to include a predicate here which can include surface expressions as arguments so that labels can be constrained in terms of other labels.
If this is done, it is assumed that the arguments are translated to the core language (see section \ref{subsec:translation-rules}) and evaluated (see section \ref{subsec:generation-rules}) before being passed to the predicate.



\subsection{The Core Language} \label{subsec:the-core-language}
\par After translation all expressions are in the core language. The core language provides better typechecking for the compiler.
\begin{figure}
\begin{center}
    \begin{align*}
        \mathit{core} &\to \mathit{core} '|_{\lab}' \mathit{core}   \\
            &\OR \mathit{core}*_{\lab}                              \\
            &\OR \mathit{core} :> \mathit{core}                     \\
            &\OR \mathit{core} :++ \mathit{core}                    \\
            &\OR (\lab \Longleftarrow b) \asgnarrow \mathit{core}   \\
            &\OR \lab                                               \\        
            &\OR <x>                                                \\
            &\OR \nil                                               \\
    \end{align*}
    \caption{The Core Language as a Context Free Grammar}
    \label{fig:the-core-language}
\end{center}
\end{figure}



\subsection{Translation Rules} \label{subsec:translation-rules}
\begin{figure*}
    % Data
     $\inference
        [AppTrans]
        {\forall arg_i \in args, }
        {<f\;args> \translate <f\;args'>}$
        % Union
     $\inference
        [UnionTrans]
        {\hastype{x}{\complex} & \hastype{y}{\complex} & x \translate x' & y \translate y'}
        {x |\{\lab\} y \translate x' \; |_{\lab} \; y'}$
     $\inference
        [UnlabelledUnionTrans]
        {\hastype{x}{\complex} & \hastype{y}{\complex} & x \translate x' & y \translate y' & \fresh{\lab'}}
        {x \; | \; y \translate (\lab' \Longleftarrow True) \mapsto x' \; |_{\lab'} \; y'}$

    % Star
     $\inference
        [\StarBaseTrans]
        {\hastype{x}{\base} & x \translate x'}
        {x*\{\lab\} \translate [x']*_{\lab}}$
     $\inference
        [\StarTrans]
        {\hastype{x}{[\complex]} & x \translate x'}
        {x*\{\lab\} \translate x'*_{\lab}}$
     $\inference
        [UnlabelledStarBaseTrans]
        {\hastype{x}{\base} & x \translate x' & \fresh{\lab'}}
        {x* \translate (\lab' \Longleftarrow True) \mapsto [x']*_{\lab'}}$
     $\inference
        [UnlabelledStarTrans]
        {\hastype{x}{[\complex]} & x \translate x' & \fresh{\lab'}}
        {x* \translate (\lab' \Longleftarrow True) \mapsto x'*_{\lab'}}$

    % Concat
     $\inference
        [\ConcatBaseTrans]
        {\hastype{x}{\base} & \hastype{y}{\base} & x \translate x' & y \translate y'}
        {x \; y \translate x' :> [y']}$
     $\inference
        [\ConcatConsTrans]
        {\hastype{x}{\complex} & \hastype{y}{[\complex]} & x \translate x' & y \translate y'}
        {x \; y \translate x' :> y'}$
     $\inference
        [\ConcatSnocTrans]
        {\hastype{x}{[\complex]} & \hastype{y}{\complex} & x \translate x' & y \translate y'}
        {x \; y \translate x' :++ [y']}$
     $\inference
        [\ConcatTrans]
        {\hastype{x}{[\complex]} & \hastype{y}{[\complex]} & x \translate x' & y \translate y'}
        {x \; y \translate x' :++ y'}$

    % Assignment
     $\inference
        [AssignmentWithPredTrans]
        {\forall arg \in args, arg \translate arg' & x \translate x'}
        {(\lab \Longleftarrow p(args, \lab)) \mapsto x \translate (\lab \Longleftarrow p(args', \lab)) \mapsto x'}$
     $\inference
        [AssignmentNoPredTrans]
        {x \translate x'}
        {\lab \mapsto x \translate (\lab \Longleftarrow True) \mapsto x'}$
    
    % Variables
     $\inference
        [Variable]
        {}
        {\lab \translate \lab}$

    % Dot
     $\inference
        [DotWithPredTrans]
        {args \; *\translate args' & fresh(\lab')}
        {(. \Longleftarrow p \; args) \translate (\lab' \Longleftarrow p(args', \lab)) \mapsto \lab'}$
     $\inference
        [\DotNoPredTrans]
        {\fresh(\lab')}
        {. \translate (\lab' \Longleftarrow True) \mapsto \lab'}$

    % Literals
     $\inference
        [LitTrans]
        {}
        {<c> \translate <c>}$ 

    % Empty
     $\inference
        [EmptyTrans]
        {}
        {\nil \translate \nil}$
\caption{\bill{Translation Rules}}
\label{fig:translation-rules}
\end{figure*}


\subsection{Dependency Map Computation} \label{subsec:dependency-map-computation}
\par the dependency map function is a bit complex and I'm not sure if we need rules for it anyways or if its sufficient to just describe the algorithm. Since we build the dependency map from a whole series of expressions representing a function I designed it so that it could be a fold function (i.e its type is dmap -> ast n a -> dmap)
\begin{enumerate}
    \item TODO: rules go here
\end{enumerate}

\subsection{Scope Checking Rules} \label{subsec:scope-checking-rules}
\par The scope checking rules aim to check if every label used in a tree is used in an appropriate place which really boils down to two rules
\begin{enumerate}
    \item No label is used without first being declared
    \item No label is used unless the dependency map allows it
\end{enumerate}
\begin{enumerate}
    \item TODO: add rules
\end{enumerate}

\subsection{Well Formedness Rules} \label{subsec:well-formedness-rules}
\par The well formedness rules aim to make sure that the depths of each expression and subexpression match the rules of our Ast n a GADT (this may already be done in the translation rules?)
\begin{enumerate}
    \item TODO: add rules
\end{enumerate}

\subsection{Generation Rules} \label{subsec:generation-rules}
\begin{figure*}
    % we assume all values are generated at assignment and therefore the environment contains them
     $\inference
        [Variable]
        {\envget{\env}{\lab} = x}
        {\inp \lab \matches x, \zout}$

     $\inference
        [\RuleEmpty]
        {}
        {\inp \nil \matches \nil, \zout}$

    % by design, the depth of r_1 and r_2 are equal
     $\inference
        [\RuleUnionLeft]
        {\envget{\env}{\lab} \mod 2 = 0 & \inpwp{\lab}{0} r_1 \matches x, \out}
        {\inp r_1 |_{\lab} r_2 \matches x, \out}$
     $\inference
        [\RuleUnionRight]
        {\envget{\env}{\lab} \mod 2 = 1 & \inpwp{\lab}{0} r_2 \matches x, \out}
        {\inp r_1 |_{\lab} r_2 \matches x, \out}$

    % Is it possible to pass the that label for the star when we switch it out for a constant label?
     $\inference
        [\RuleStarInit]
        {\envhasnt{\pathvar}{\lab} & \envget{\env}{\lab} = n & \inpwp{\lab}{0} r*_{\lab} \matches xs, \out}
        {\inp r*_{\lab} \matches xs, \out}$
     $\inference
        [\RuleStar]
        {\envget{\pathvar}{\lab} < n & \envget{\env}{\lab} = n & \inp r \matches xs, \out & \inpwp{\lab}{\envget{\pathvar}{\lab} + 1} r*_{\lab} \matches ys}
        {\inp r*_{\lab} \matches xs ++ ys}$
     $\inference
        [\RuleStarEnd]
        {\envget{\env}{\lab} \leq 0 \lor \envget{\pathvar}{\lab} \geq \envget{\env}{\lab}}
        {\inp r*_{\lab} \matches \nil, \out}$

    $\inference
        [\RuleStar]
        { \envget{\env}{\lab} = n
        & xs = \forall 0 \leq i < n . \inpwp{\lab}{i} r \matches xs_i, \out
        }
        {\inp r*_{\lab} \matches \foldconcat{xs}, \out}$

    % by design, r_1 and r_2 have a non-zero depth for concat
     $\inference
        [\RuleConcat]
        {\inp r_1 \matches xs, \out & r_2 \matches ys, \dout}
        {\inp r_{1} :++ r_{2} \matches xs ++ ys, \dout}$
     $\inference
        [\RuleCons]
        {\inp r_1 \matches x, \out & \out r_1 \matches xs, \dout}
        {\inp r_{1} :> r_{2} \matches x : xs, \dout}$

    % TODO be more clear about difference between this and apply maybe?
     $\inference
        [\RuleLiteral]
        {}
        {\inp <c> \matches c, \out}$

     $\inference
        [\RuleApply]
        {\inp args \matches args' \env}
        {\inp <f \; args> \matches f(args'), \env}$

     $\inference
        [\RuleAssign]
        {\arbitrary{x} & p(args, x)}
        {\inp (\lab \Longleftarrow p(args, \lab)) \mapsto e, \outw{\lab}{x}}$
    \caption{\bill{Generation Rules}}
    \label{fig:generation-rules}
\end{figure*}


\par The generation rules demonstrate how a nested list is generated from a symbolic expression. The rules are constructed such that you are guaranteed that the nested depth of the generated value will be the same as the depth index of the core AST that generated it.


\subsection{Check Rules} \label{subsec:check-rules}
\par The check rules aim to determine if it is possible that a given nested list could be generated by a given ast. I'm gonna need some special syntax here that means "there exists a prefix of this list such that..." and "and the rest of the list has some other property" for concat, star, and maybe wrap/cons
\begin{enumerate}
    \item TODO: add rules
\end{enumerate}

